\documentclass[aspectratio=169,dvipdfmx,12pt,a4j,handout]{beamer}
\usepackage{bxdpx-beamer}
\usepackage{pgfpages}
\usepackage{pxjahyper}
\usepackage{graphicx}
\pgfpagesuselayout{8 on 1}[a4paper,landscape,border shrink=5mm]
\pgfpageslogicalpageoptions{1}{border code=\pgfusepath{stroke}}
\pgfpageslogicalpageoptions{2}{border code=\pgfusepath{stroke}}
\pgfpageslogicalpageoptions{3}{border code=\pgfusepath{stroke}}
\pgfpageslogicalpageoptions{4}{border code=\pgfusepath{stroke}}
\pgfpageslogicalpageoptions{5}{border code=\pgfusepath{stroke}}
\pgfpageslogicalpageoptions{6}{border code=\pgfusepath{stroke}}
\pgfpageslogicalpageoptions{7}{border code=\pgfusepath{stroke}}
\pgfpageslogicalpageoptions{8}{border code=\pgfusepath{stroke}}
\usetheme{AnnArbor}
\renewcommand{\kanjifamilydefault}{\gtdefault}


\title{Beamer課題}
\subtitle{コンピュータリテラシ}
\author[神野]{神野 友哉 2022311042}
\institute[愛知県立大]{\Large 愛知県立大学　情報科学部}
\date{\large \today}


\begin{document}

\begin{frame}
  \titlepage
\end{frame}

\begin{frame}\frametitle{プレゼンの流れ}
  \tableofcontents
\end{frame}

\section{特技}
\begin{frame}\frametitle{特技}
 私の特技は、短時間睡眠でもある程度頭が回ることです。\pause
この生来の特技のおかげで、ほぼすべての講義に間に合うことが可能です。\pause
また、辛いものに対する耐性もあります。
\begin{figure}[htb]
  \centering
  \includegraphics[width=0.4\linewidth]{suimin_man.png}
\end{figure}
\end{frame}

\section{過去最高のエピソード}
\begin{frame}\frametitle{過去最高のエピソード}
私のこの大学に入学して以来の、過去最高のエピソードは、
夜の８時に家に帰宅したことです。\pause
今までの学生人生で、ここまで遅くまで課題をしていたことはありませんでした。\pause
大学生になったと感じます。
\begin{figure}[htb]
  \centering
  \includegraphics[width=0.2\linewidth]{daigakusei_man.png}
\end{figure}
\end{frame}

\section{大学数学で習得した公式}
\begin{frame}\frametitle{大学数学で習得した公式}
私は、大学に入学して以来、多くの数学の公式を会得しましたが、\pause 特に私の記憶に残っている公式は、以下の式で表されるような、マクローリンの定理です。\pause
\vskip2\baselineskip
fは[0,x]上でn−1回連続微分可能。(0,x)上でn回連続微分可能とする。このとき
\[
 f(x)=f(0)+\frac{f(0)}{1!}x+\frac{f^2(0)}{2!}x^2+\cdots +\frac{f^n(c)}{n!}x^n
\]
を満たすcが0とxの間に存在する。\pause
\vskip2\baselineskip
という内容です。
\end{frame}

\subsection{マクローリン展開の利用法}
\begin{frame}\frametitle{マクローリン展開の利用法}
 円周率の値や様々な三角関数の値，対数の値など，マクローリン展開によって詳しい近似値を求められます。\pause これにより携帯や3Dなど、今日における情報社会の地盤が形作られているといっても過言ではありません。
\end{frame}

\section{理系の文章術}
\begin{frame}\frametitle{理系の文章術からの知見}
  私は理系の文章術という、更科功(さらしなこう)氏著書の文章本を入手しました。\pause
  印象に残った部分は、"わかりやすい文章は最高の文章ではない。\pause 'わかりやすい'と'正確さ'には負の相関があり、論文などの文章では、妥当性が失われる場合がある。\pause しかし、'正確さ'を損なわない範囲で、わかりやすくすることは重要である。"という内容でした。
\end{frame}

\subsection{その他の知見}
\begin{frame}\frametitle{その他の知見}
  あまり意識しない句読点の打ち方について、\pause ”文を中止する時、\pause 副詞(句)の前後、\pause 動詞を受ける言葉が遠い時の受ける言葉の後、\pause 意味のまとまりの前後に打つ”\pause と具体的なルールを学ぶことができました。
\end{frame}

\end{document}
